\section{Joint response and an exact finite-dimensional remainder}
\label{sec:joint}
The compact error in Section~\ref{sec:feedback} can be replaced by a
finite-rank error. The construction uses a proved positive high
sector and keeps every coupling to its finite-dimensional
complement. It leads to an exact finite matrix criterion, whose
sign remains to be established.

\subsection{Choose the high sector without the Weil spectrum}
For $q\in(0,1)$ put
\[
r_n(q)=\max_{1\le i\le n}
 \left(\sum_{j=1}^{i-1}q^j+\sum_{j=1}^{n-i}q^j\right),\qquad
n_p=\lceil L/\log p\rceil.
\]
The row-sum bound on the finite prime chains in
\eqref{prime:finite-contact}, the negative odd pole bound and
Theorem~\ref{source:bracket} give
\begin{equation}\label{joint:beta}
W_L\ge B(H_{\rm N})-\beta_LI,\qquad
\beta_L=-w_0+2\sinh(L/2)-L+
 \sum_{\log p<L}(\log p)r_{n_p}(p^{-1/2}).
\end{equation}
The coarser replacement of the last sum by
$2\sum_{m\log p<L}(\log p)p^{-m/2}$ is also valid.

Use the explicit Neumann modes
\[
e_0=L^{-1/2},\qquad
e_j(x)=\sqrt{2/L}\cos(j\pi(x+L/2)/L)\quad(j\ge1),
\qquad P=P_N=\sum_{j=0}^{N-1}e_je_j^*,\quad Q=I-P.
\]
Choose $N\ge1,\delta>0$ such that
\begin{equation}\label{joint:cutoff}
B((N\pi/L)^2)\ge\beta_L+\delta.
\end{equation}
Such choices exist at each fixed $L$. The restricted closed
arithmetic form then satisfies $W_L|_{QX_L}\ge\delta I$.
This is not a statement about the mixed low/high terms.
The vectors $e_j$ belong even to $\Dom T_L$: their zero-extension
Fourier transforms are $O_j((1+|\tau|)^{-1})$, which remain
square-integrable after multiplication by $B(\tau^2)$.
Thus $P$ preserves the form domain and $QW_LP$ is bounded and
finite rank.

There is a wholly rational example at $L=1$. Only prime two's
first return is active, and
\[
\beta_1<
\frac{58}{100}+\frac{11}{7}+\frac{21}{10}+\frac{23}{20}
       +\frac{43}{1000}+\frac12
=\frac{41611}{7000}.
\]
The bounds used here are
$\gamma_{\!E}<0.58$, $\pi<22/7$, $\log2<0.7$,
$\log\pi<23/20$, $2\sinh(1/2)-1<0.043$, and
$\log2/\sqrt2<1/2$. They follow from positive or alternating
rational series. For example
$\log2=2\sum_{j\ge0}3^{-(2j+1)}/(2j+1)>6931/10000$,
so $\gamma_{\!E}<H_{256}-8(6931/10000)<58/100$.
The integral identities for $\arctan1$ and
$22/7-\pi=\int_0^1x^4(1-x)^4/(1+x^2)\,dx$ give $3<\pi<22/7$.
A finite exponential sum bounds $e^{23/20}$ from below by $22/7$;
a geometric tail bounds the positive sinh series.
Finally, exact fraction arithmetic gives
\[
B((4\pi)^2)>
\sum_{k=0}^{63}\frac2{a_k}\frac{144}{a_k^2+144}
>\frac{3007}{500}>
\frac{41611}{7000}+\frac1{16}.
\]
Thus $N=4$, $\delta=1/16$ is a proved high-mode cutoff.
It is not a certificate for the full form at $L=1$.

\subsection{One field drives all channels}
Combine the already specified sources in the same positive space:
\begin{equation}\label{joint:reference}
A=(A_0,J),\qquad
T=A^*A=T_0+J^*J=W_L+\alpha I,\qquad
\alpha=e^D+c_{\mathcal S}-w_0=2\kappa,\qquad T\ge mI.
\end{equation}
The domain is $\mathcal D_{\log,L}$ and $\Dom T=\Dom T_L$.
Prepare a common field in both the gamma-and-pole source and
every prime graph by minimizing
$\norm{A_0u}^2+\norm{Ju}^2-2\re\ip v u$.
Its response is $T^{-1}v$. A first joint readout obeys
\[
\norm{A(I-\alpha T^{-1}/2)f}^2
 =Q_L[f]+\frac{\alpha^2}{4}\ip f{T^{-1}f}.
\]
Since $T\ge T_0>0$, this error is no greater than the earlier
massless error. The full repeated response gives a stronger result.

Restrict $A$ to $Q\mathcal D_{\log,L}$, denoting it $A_H$, and put
\[
T_H=A_H^*A_H,\qquad V=T_H^{-1},\qquad
H=T_H-\alpha I\ge\delta I,\qquad
B=QTP=QW_LP,\qquad M=PTP.
\]
In this section $B$ without an argument is the finite coupling
matrix; $B(s)$ remains the gamma multiplier.
The arithmetic content of the coupling is explicit:
\[
B=QT_LP+QP_LP-
\sum_{m\log p<L}(\log p)p^{-m/2}
                         Q(T_{m\log p}+T_{m\log p}^*)P .
\]
For $f=p+q$, with $p=Pf$, $q=Qf$, the exact form is
\begin{equation}\label{joint:blocks}
Q_L[p+q]=\ip q{Hq}+2\re\ip q{Bp}
                          +\ip p{(M-\alpha I)p}.
\end{equation}
Terms involving $H$ are interpreted as forms on its form domain.

\subsection{A convergent response series with all mixed returns}
Let
\[
a_j=\frac{\binom{2j}{j}}{4^j(2j-1)},\qquad
a_1=\frac12,\quad a_{j+1}=a_j\frac{2j-1}{2j+2},\qquad
\theta=\frac{\alpha}{\alpha+\delta}<1.
\]
The operator series
\begin{equation}\label{joint:series}
R=I-\sum_{j\ge1}a_j\alpha^jV^j,\qquad
R_\ell=I-\sum_{j=1}^{\ell}a_j\alpha^jV^j
\end{equation}
converges in norm and gives a positive invertible $R$ with
\[
R^2=I-\alpha V,\qquad RT_HR=H,\qquad
\norm{A_HRq}^2=\ip q{Hq}.
\]
It is precisely the positive high-sector factor
$(I-\alpha V)^{1/2}$. Its use is justified by the independent
bound \eqref{joint:cutoff}; no sign of the full $W_L$ has been
assumed.
The stronger estimate needed for the unbounded output is
\begin{equation}\label{joint:tail}
\norm{A_H(R_\ell-R)}
\le\frac{a_{\ell+1}\alpha^{\ell+1}}
 {(1-\theta)(\alpha+\delta)^{\ell+1/2}}\longrightarrow0.
\end{equation}
Indeed the tail on an eigenvalue $\lambda\ge\alpha+\delta$ is
$\sqrt\lambda\sum_{j>\ell}a_j(\alpha/\lambda)^j$.
Moreover
\begin{equation}\label{joint:positive-error}
E_\ell:=T_HR_\ell^2-H\ge0,\qquad
\norm{E_\ell}\le
\frac{2a_{\ell+1}\alpha\theta^\ell}{1-\theta},
\quad\ell\ge1.
\end{equation}
Each $E_\ell$ is compact. A small ratio $\delta/\alpha$ can make
this convergence slow.

To specify the noncommuting returns, write
$T_H=T_{0,H}+K_H$ with $K_H=(J|_Q)^*(J|_Q)\ge0$.
Choose the known bound
$t=\sum_{p\in\mathcal S}4d_pq_p/(1-q_p^2)\ge\norm{K_H}$ and put
$V_0=(T_{0,H}+tI)^{-1}$. Then
\begin{equation}\label{joint:words}
V=\sum_{j\ge0}[V_0(tI-K_H)]^jV_0,\qquad
\norm{V_0(tI-K_H)}\le\frac{t}{m+t}<1,
\end{equation}
with just $V_0$ when $t=0$. Here
\[
K_H=c_{\mathcal S}I_{QX_L}-
\sum_{m\log p<L}(\log p)p^{-m/2}
                      Q(T_{m\log p}+T_{m\log p}^*)Q.
\]
The order of the intervening resolvents and compressions in
\eqref{joint:words} retains every mixed gamma/prime return.
Replacing these ordered words by scalar Euler multipliers would
change the operator. Also
$H^{-1}=\sum_{j\ge0}\alpha^jV^{j+1}$ converges in norm, supplying
the second high-sector response below.

\subsection{Retain the unused output and obtain finite-rank error}
The range of $A_H$ is closed. Its orthogonal projection is the
bounded extension $\Pi_H=A_HVA_H^*$. On a low input,
\[
Ap=A_HVBp+Zp,\qquad
Z=(I-\Pi_H)A|_{PX_L},\qquad
G:=Z^*Z=M-B^*VB\ge mI.
\]
The last bound follows by minimizing the known positive $T$-energy
over the high input at fixed $p$. In particular no unknown
arithmetic Schur complement has been declared positive.

\begin{theorem}\label{joint:completion}
In the original positive output space, define
\begin{equation}\label{joint:source}
\Phi_L(p+q)=A_HR(q+H^{-1}Bp)+Zp .
\end{equation}
It is a closed source on $\mathcal D_{\log,L}$, and for all
inputs in that domain its polarized Gram identity is
\begin{equation}\label{joint:gram}
\ip{\Phi_Lf}{\Phi_Lg}=Q_L(f,g)+\ip{Pf}{D_NPg},
\qquad
D_N=\alpha I+B^*(H^{-1}-V)B
    =\alpha(I+B^*VH^{-1}B)>0.
\end{equation}
The error operator $PD_NP$ has rank exactly $N$.
\end{theorem}
\begin{proof}
The two outputs in \eqref{joint:source} are orthogonal. Expanding
their norm with $RT_HR=H$ gives
$\ip q{Hq}+2\re\ip q{Bp}
+\ip p{(G+B^*H^{-1}B)p}$.
Subtract \eqref{joint:blocks}, use $G=M-B^*VB$, and use
$H^{-1}-V=\alpha VH^{-1}\ge0$ to obtain \eqref{joint:gram}.
Polarization gives the bilinear identity.
For closedness subtract the original source:
\begin{equation}\label{joint:correction}
\Phi_L-A=A_H(R-I)Q+
                 A_H(RH^{-1}-V)BP=:F.
\end{equation}
The first term extends to a compact bounded operator since its
singular magnitude at $\lambda\in\operatorname{spec}T_H$ is
\[
\sqrt\lambda(1-\sqrt{1-\alpha/\lambda})
 =\frac{\alpha}{\sqrt\lambda(1+\sqrt{1-\alpha/\lambda})}
 \longrightarrow0.
\]
The second is bounded and finite rank. Thus $\Phi_L=A+F$ is a
bounded perturbation of a closed source.
\end{proof}

It has the same graph topology and smooth compact-support core
as $A$, and its maximal adjoint is
$\Phi_L^*=A^*+F^*$ on $\Dom A^*$.
Consequently
$\Phi_L^*\Phi_L=W_L+PD_NP$ on $\Dom T_L$.
Using $R_\ell$ instead of $R$ gives the exact additional error
\[
(Q+H^{-1}BP)^*E_\ell(Q+H^{-1}BP)\ge0,
\]
which tends to zero in operator norm by
\eqref{joint:positive-error}.
Finite series generally leave an infinite-rank extra error.

The correction $F$ itself is not finite rank. Bounded-perturbation
min--max and finite-codimension interlacing give
$\lambda_n(T_H)=\log n+O_L(1)$.
The displayed singular magnitude and finite-rank interlacing in
\eqref{joint:correction} imply
\begin{equation}\label{joint:singular}
\sqrt{\log n}\,s_n(F)\longrightarrow\alpha/2.
\end{equation}
Thus the source correction remains outside every finite Schatten
class, consistently with Section~\ref{sec:feedback}. Finite rank
of the form error does not imply finite rank of the source change.

\subsection{The exact remaining matrix criterion}
Completing the high square in \eqref{joint:blocks} gives
\begin{equation}\label{joint:schur}
Q_L[p+q]=\norm{A_HR(q+H^{-1}Bp)}^2+\ip p{S_Np},
\qquad
S_N=M-\alpha I-B^*H^{-1}B=G-D_N.
\end{equation}
Since $q=-H^{-1}Bp$ is admissible for every $p$,
\begin{equation}\label{joint:criterion}
Q_L\ge0\quad\Longleftrightarrow\quad S_N\ge0
\quad\Longleftrightarrow\quad G^{-1/2}D_NG^{-1/2}\le I.
\end{equation}
Both $G$ and $D_N$ are known positive matrices; their difference
has no established sign. At $L=1$ the rational cutoff above makes
this a four-by-four comparison, but no entries or positivity
certificate for that arithmetic matrix are asserted.

A rational non-arithmetic control makes the logical limit explicit.
Take $\alpha=1$,
\[
T_H=\operatorname{diag}(25/16,25/9),\quad
H=\operatorname{diag}(9/16,16/9),\quad
B=\begin{pmatrix}1&2\\3&1\end{pmatrix},\quad
Z_c=\begin{pmatrix}2&1\\0&1\end{pmatrix}.
\]
On two low and two high coordinates use the invertible source
\[
A=\begin{pmatrix}T_H^{-1/2}B&T_H^{1/2}\\Z_c&0\end{pmatrix},
\qquad
\Phi=\begin{pmatrix}H^{-1/2}B&H^{1/2}\\Z_c&0\end{pmatrix}.
\]
Their Gram identity holds exactly, with
\[
D_N=\frac1{3600}
\begin{pmatrix}14257&10379\\10379&20713\end{pmatrix},\qquad
S_N=\frac1{3600}
\begin{pmatrix}143&-3179\\-3179&-13513\end{pmatrix}.
\]
The vector $p=(0,1)$, $q=(-32/9,-9/16)$ has
$\ip{p+q}{(A^*A-I)(p+q)}=-13513/3600<0$.
Thus a positive shifted source, a positive high sector and the
completion identity can coexist with a negative full target.
The reduction is valid for a wider class of semibounded forms;
its existence alone cannot be an arithmetic positivity proof.

\subsection{Compatibility inside a fixed support envelope}
Construct $\Phi_*$ once on $I_{L_*}$, fixing its delay, prime set
and mode projection. For $L\le L_*$ let $j_L$ be zero extension
into that envelope. The restricted sources are closed on the
full logarithmic domains, have the same smooth cores, and obey
\begin{equation}\label{joint:envelope}
\norm{\Phi_*j_Lf}^2
=Q_L[f]+\ip f{j_L^*P_*D_*P_*j_Lf}.
\end{equation}
The complete arithmetic form is unchanged under the extension;
the additional inactive prime overlaps vanish. These sources
are compatible with nested intervals inside the fixed envelope.
Rebuilding the source on a larger envelope generally changes
the compressed error. No single all-support source or all-prime
limit, and no rule allowing these errors to be subtracted, has
been obtained.

\subsection{A bound on the mixed response with its residual retained}
The boundary identity in Theorem~\ref{boundary:operator} makes the
gamma part of the mixed block explicit. Since $P$ commutes with
$B(H_{\rm N})$,
\begin{equation}\label{mixed:boundary}
B=Q\mathcal K_LP+QP_LP-QJ_L^{\rm off}P.
\end{equation}
Here $P_L$ denotes the pole operator and $J_L^{\rm off}$ the complete
active prime translation sum. All three terms and their interference
are retained in the following estimate.

Define the known positive diagonal comparison on $QX_L$ by
\[
\Lambda_Ne_j=d_je_j,\qquad
d_j=B((j\pi/L)^2)-\beta_L\ge\delta\quad(j\ge N).
\]
Then $H\ge\Lambda_N$ as forms.
\begin{proposition}\label{mixed:residual}
For any map $Y:PX_L\to\Dom H$, set
$R_Y=B-HY$ and $U_Y=B^*Y+Y^*B-Y^*HY$. Then
\begin{equation}\label{mixed:identity}
B^*H^{-1}B=U_Y+R_Y^*H^{-1}R_Y,
\qquad
U_Y\le B^*H^{-1}B\le U_Y+R_Y^*\Lambda_N^{-1}R_Y.
\end{equation}
For $K>N$, let $P_{N,K}$ project onto modes $N,\ldots,K-1$ and
let $Y_K$ solve
\[
(P_{N,K}HP_{N,K})Y_K=P_{N,K}B,
\qquad Y_K:PX_L\longrightarrow P_{N,K}X_L.
\]
This high compression is independently positive. Its residual
$R_K=B-HY_K$ is supported on modes $j\ge K$, and
\begin{equation}\label{mixed:certificate}
S_N\ge PW_LP-B^*Y_K-d_K^{-1}R_K^*R_K.
\end{equation}
\end{proposition}
\begin{proof}
Expand the residual square without reordering factors, and use
$H^{-1}\le\Lambda_N^{-1}$. Galerkin orthogonality gives
$P_{N,K}R_K=0$ and
$Y_K^*HY_K=Y_K^*B=B^*Y_K$. On the residual's remaining cosine
support, $\Lambda_N^{-1}\le d_K^{-1}I$.
\end{proof}
This estimates the infinite response through a full residual Gram.
For a residual with rows $r_j=e_j^*R$, its diagonal inverse tail
also satisfies the sharper bound
\begin{equation}\label{mixed:mode-tail}
R^*\Lambda_N^{-1}R\le
\sum_{j=N}^{K-1}\frac{r_j^*r_j}{d_j}
+\frac1{d_K}\left(R^*R-\sum_{j=N}^{K-1}r_j^*r_j\right).
\end{equation}
If $\widetilde R$ is an approximation with
$\norm{\Lambda_N^{-1/2}(R-\widetilde R)}\le\epsilon$, the additional
upper error in this Gram is at most
$(2\norm{\Lambda_N^{-1/2}\widetilde R}\epsilon+\epsilon^2)I$.
For finite cosine trial columns, \eqref{boundary:column-tail}
controls mass truncation in both $B$ and $HY$.
The full Gram of the finite mass, pole and translated cosine
columns can be bounded by elementary integrals; omitted mode
rows are then retained through \eqref{mixed:mode-tail}.

At fixed $L,N$, Galerkin convergence in the form core gives
$Y_K\to H^{-1}B$. Write $H=\Lambda_N+E_N$ with $E_N$ bounded,
using Theorem~\ref{boundary:operator}. Then
$R_K=(I-P_K)(B-E_NY_K)\to0$ in norm on the finite input space.
Thus the right side of \eqref{mixed:certificate} converges to $S_N$.
If $S_N$ is strictly positive at this length, some finite trial
will certify it with sufficiently accurate enclosures. No such
eventual certificate is asserted on a null direction when $S_N$
is only semidefinite.
The remaining arithmetic lemma is a choice of $K(L)$ and a proof
that the right side of \eqref{mixed:certificate}, including all
enclosure errors, is nonnegative for every required $L$.
The residual theorem does not assert that sign.

The same distinction applies to ordered inverse truncations.
With $F_\ell=\sum_{j=0}^{\ell}\alpha^jV^{j+1}$,
\[
0\le H^{-1}-F_\ell\le\theta^{\ell+1}\delta^{-1}I.
\]
Consequently $PW_LP-B^*F_\ell B$ is an \emph{upper} bound for $S_N$.
Its positivity is insufficient. For any $\ell\ge0$, the rational
control
\[
W=\begin{pmatrix}a&1\\1&1\end{pmatrix},\quad
a=1-2^{-\ell-2},\quad \alpha=1
\]
has $W+I>0$, $H=1$, $V=1/2$, but
$a-F_\ell=2^{-\ell-2}>0$ and $S=a-1=-2^{-\ell-2}<0$.
The missing inverse tail reverses the apparent conclusion.
